\chapter{Regular coefficient models}
\label{v4-ch:regular-coefficient-models}

An associative coefficient algebra already produces fields:
translate the first input and multiply it by the second.
Their products have no poles. Nevertheless, their realization
on collision strata must retain the order of the inputs.
For the matrix units
\[
 a=c=E_{12},\qquad b=E_{21},
 \qquad ac=0,\qquad abc=E_{12},
\]
combining the outer coefficients before allowing the middle one
to act changes a nonzero operation into zero. The obstruction
concerns a partial collision, where the first and third positions
agree and the second remains separate.

We construct coefficients on every such stratum by retaining each
consecutive run of a position in the original word. Their
principal-parts maps form a complex, and their collision and cut
maps are coherent. A specified presentation permits recovery of
the original differential algebra. We then construct a different
coefficient system using an explicit exchange. It compresses the
runs while retaining their evaluated products. The compression
has a computable kernel; it need not be an equivalence.

\section{Regular fields and the information in a word}
\label{v4-sec:regular-fields}

Fix a field \(k\) of characteristic zero and \(X=\mathbb A^1_k\)
with coordinate \(z\). A \emph{differential coefficient algebra}
is a unital differential graded associative \(k\)-algebra
\((B,d,\mu,1)\), with \(1\ne0\), together with a degree-zero
derivation \(T\) commuting with \(d\). No finite-dimensionality
or boundedness condition is imposed on \(B\). Tensor products
in this chapter are algebraic over \(k\), unless another
coefficient category is specified.

Put
\begin{equation}\label{v4-eq:regular-field}
 Y_B(a,u)b=(e^{uT}a)b
          =\sum_{r\geq0}\frac{u^r}{r!}(T^ra)b
          \quad\text{in }B[[u]].
\end{equation}
This is a formal power series, even when \(T\) is not locally
nilpotent. Its coefficient in every fixed degree is defined.
Writing \(Y_B(a,u)b=\sum_{n\in\mathbb Z}a_{(n)}b\,u^{-n-1}\),
we obtain
\[
 a_{(-r-1)}b=\frac{(T^ra)b}{r!}\quad(r\geq0),
 \qquad a_{(n)}b=0\quad(n\geq0).
\]
In particular, \(a_{(-1)}b=ab\).

\begin{proposition}\label{v4-prop:regular-field-identities}
The fields \eqref{v4-eq:regular-field} preserve the differential
and satisfy the vacuum and translation identities
\[
 Y_B(1,u)b=b,\qquad Y_B(a,u)1=e^{uT}a,
 \qquad Y_B(Ta,u)b=\partial_uY_B(a,u)b,
\]
\[
 T(Y_B(a,u)b)=Y_B(Ta,u)b+Y_B(a,u)Tb.
\]
Their simultaneous product, rooted at the last position, is
\begin{equation}\label{v4-eq:regular-simultaneous}
 m_n(b_1,\ldots,b_n;z)
       =\prod_{i=1}^n e^{(z_i-z_n)T}b_i
\end{equation}
in the original order. Consecutive substitutions give this
same power series.
\end{proposition}

\begin{proof}
The unit and derivation identities give the displayed formulas
coefficient by coefficient. For example,
\[
 \begin{split}
 Y_B(Y_B(a,u)b,v)c
   &=e^{vT}\bigl((e^{uT}a)b\bigr)c\\
   &=(e^{(u+v)T}a)(e^{vT}b)c
     =Y_B(a,u+v)Y_B(b,v)c.
 \end{split}
\]
There are no negative powers to obstruct the substitution.
At several consecutive blocks, the same derivation identity
moves each block to the common root. Associativity removes
the parentheses. The graded Leibniz rule and \(dT=Td\)
prove compatibility with the differential.
\end{proof}

The order in \eqref{v4-eq:regular-simultaneous} remains part
of the operation. For constant fields, changing it replaces
\(ab\) by \(ba\). A rule relating those expressions requires
additional maps.

To measure the information between occurrences, leave \(r\)
inputs open and define
\begin{equation}\label{v4-eq:regular-open-slots}
 \Gamma_{B,r}:B^{\otimes(r+1)}
       \longrightarrow\operatorname{Hom}_k(B^{\otimes r},B),
 \qquad
 \Gamma_{B,r}(a_0,\ldots,a_r)(x_1,\ldots,x_r)
       =a_0x_1a_1\cdots x_ra_r .
\end{equation}

\begin{proposition}[Matrix coefficients with open inputs]
\label{v4-prop:regular-matrix-context}
For \(B=\operatorname{Mat}_d(k)\), the map
\(\Gamma_{B,r}\) is an isomorphism. A linear encoding of its
source that retains every evaluation therefore has dimension
at least \(d^{2r+2}\). For \(r=1\) and \(d>1\), one copy of
\(B\) cannot supply that encoding.
\end{proposition}

\begin{proof}
On matrix units, the value is
\[
 \Gamma_{B,r}(E_{i_0j_0},\ldots,E_{i_rj_r})
       (E_{p_1q_1},\ldots,E_{p_rq_r})
 =\left(\prod_{\ell=1}^r
        \delta_{j_{\ell-1},p_\ell}\delta_{q_\ell,i_\ell}\right)
         E_{i_0j_r}.
\]
Every elementary linear map between the indicated tensor bases
occurs exactly once. This proves the isomorphism. If
\(\Gamma_{B,r}\) factors as \(B^{\otimes(r+1)}\to L\to
\operatorname{Hom}_k(B^{\otimes r},B)\), its first map is
injective, which proves the dimension bound.
\end{proof}

There is a stronger constraint when the encoding is required
to be natural in the coefficient algebra.

\begin{proposition}[Naturality retains the open word]
\label{v4-prop:regular-natural-context}
Suppose a quotient
\[
 q_B:B^{\otimes(r+1)}\twoheadrightarrow Q_B
\]
is given for every differential coefficient algebra, is natural
under injective unital maps commuting with \(d,T\), and permits
\(\Gamma_{B,r}\) to factor through it. Then every \(q_B\) is
an isomorphism. For finite-dimensional ordinary algebras, it
suffices to require naturality under injections into other
finite-dimensional algebras.
\end{proposition}

\begin{proof}
Adjoin noncommuting symbols \(X_1,\ldots,X_r\) with
\(dX_i=TX_i=0\). The free \(B\)-algebra on these symbols is
the tensor algebra over \(B\) of
\(B\otimes_k k\{X_1,\ldots,X_r\}\otimes_k B\).
Its summand with symbol word \(X_1\cdots X_r\) is a copy
of \(B^{\otimes(r+1)}\). Thus
\[
 a_0\otimes\cdots\otimes a_r
       \longmapsto a_0X_1a_1\cdots X_ra_r
\]
is injective. If \(w\in\ker q_B\), naturality sends \(w\)
into the kernel of the corresponding quotient for the free
algebra. Evaluating its open inputs at \(X_1,\ldots,X_r\)
therefore gives zero. Injectivity implies \(w=0\).
For finite-dimensional \(B\), quotient the free algebra by
symbol words of length greater than \(r\). This algebra is
finite-dimensional and retains both the injection of \(B\)
and the separating summand.
\end{proof}

At one fixed algebra the kernel can be nonzero.
For \(B=k[\epsilon]/(\epsilon^2)\),
\(\epsilon\otimes1-1\otimes\epsilon\) evaluates to zero
on every element of \(B\). Under
\(\epsilon\mapsto E_{12}\), its evaluation on \(E_{21}\)
becomes \(E_{11}-E_{22}\). Enlarging the allowed intermediate
inputs reveals information that the original evaluations lost.

\section{Run coefficients and their residue maps}
\label{v4-sec:regular-runs}

Use the word category \(\mathcal W_X\) and its cut product
from Chapter~\ref{v4-ch:words-collisions-units}.
For a partition \(\pi\) of an ordered set \(I\), read its
equality blocks along the word and remove only adjacent
repetitions. The resulting list is the list of \emph{runs}.
Write \(\ell(\pi)\) for its length and \(\beta_\pi(r)\)
for the block occupied by run \(r\). At three positions,
\[
 \begin{array}{c|c|c}
 \pi&\text{run list}&\text{coefficient}\\ \hline
 12|3&(12,3)&B^{\otimes2}\\
 1|23&(1,23)&B^{\otimes2}\\
 13|2&(13,2,13)&B^{\otimes3}.
 \end{array}
\]
The equality \(z_1=z_3\) does not make those two occurrences
adjacent in the original word.

A coarsening \(\pi\leq\rho\) relabels every run by its
\(\rho\)-block. Multiply the factors in each newly consecutive
run, retaining their order. This defines
\begin{equation}\label{v4-eq:regular-run-map}
 p_{\pi\rho}:B^{\otimes\ell(\pi)}
       \longrightarrow B^{\otimes\ell(\rho)}.
\end{equation}
For example, \(p_{1|2|3,13|2}\) is the identity and
\(p_{13|2,123}(a,b,c)=abc\).

\begin{lemma}\label{v4-lem:regular-run-functor}
The maps \eqref{v4-eq:regular-run-map} are chain maps and obey
\[
 p_{\rho\sigma}p_{\pi\rho}=p_{\pi\sigma}.
\]
For each block \(C\) of \(\pi\), put
\(T_C=\sum_{\beta_\pi(r)=C}T^{(r)}\).
Along a merger, \(p_{\pi\rho}\) intertwines the sum of
the input block translations with the output block translation.
\end{lemma}

\begin{proof}
A final run is an interval of the original word. An intermediate
coarsening merely parenthesizes the ordered product on that
interval. Associativity gives the composition identity.
The two derivation identities give the assertions for \(d\)
and \(T\). No coefficient factors are permuted in this argument.
\end{proof}

Let \(X^\pi\) have one coordinate per equality block,
ordered by first occurrence. Let
\[
 i_\pi:X^\pi\hookrightarrow X^I,\qquad
 j_\pi:\operatorname{Conf}_\pi(X)\hookrightarrow X^\pi
\]
be the diagonal embedding and the open subset where different
blocks occupy distinct positions. On \(X^\pi\), set
\[
 M_\pi(B)=\omega_{X^\pi}\otimes_k B^{\otimes\ell(\pi)}.
\]
The right differential-module structure is
\begin{equation}\label{v4-eq:regular-run-connection}
 (f\,\omega_\pi\otimes\mathbf b)\partial_{z_C}
       =-\partial_{z_C}f\,\omega_\pi\otimes\mathbf b
        -f\,\omega_\pi\otimes T_C\mathbf b,
 \qquad
 \omega_\pi=dz_{C_1}\wedge\cdots\wedge dz_{C_{|\pi|}}.
\end{equation}
Define the graded object
\begin{equation}\label{v4-eq:regular-run-complex}
 C_I^{\mathrm{run}}(B)
    =\bigoplus_{\pi\in\Pi(I)}
       i_{\pi,*}j_{\pi,*}j_\pi^*M_\pi(B)[|\pi|],
 \qquad C_\varnothing^{\mathrm{run}}(B)=k.
\end{equation}
The sum is finite for each \(I\). A coefficient of internal
degree \(q\) in its \(r\)-block term has total degree \(q-r\).

We specify the boundary, including its normal module.
For a normal coordinate \(u\), let \(e_j\), \(j\geq0\),
denote the class of \(u^{-j-1}\) times its normal density.
In right-module notation,
\[
 e_0u=0,\qquad e_ju=e_{j-1}\ (j>0),\qquad
 e_j\partial_u=(j+1)e_{j+1}.
\]
Writing \(\delta=e_0\), equivalently
\(e_j=\delta\partial_u^j/j!\).
The map \(\operatorname{pp}_u\) takes the negative part
of a Laurent expansion to this full normal module. It retains
all negative powers.

For an elementary merger of blocks \(i<j\), take
\(u=z_{C_i}-z_{C_j}\), retain \(z_{C_j}\) as the merged
coordinate, and put the merged block in position \(i\).
The face on unshifted coefficients is
\begin{equation}\label{v4-eq:regular-run-face}
 D_{\pi\rho}(f\otimes\mathbf b)
       =(-1)^{j-2}\operatorname{pp}_u
          \bigl(f\,p_{\pi\rho}e^{uT_{C_i}}\mathbf b\bigr).
\end{equation}
All other differences are expanded normally before taking
principal parts. The sign is fixed by
\(\omega_\pi=(-1)^{j-2}du\wedge\omega_\rho\).
Normal transfer extends this formula to the indicated
diagonal direct image. A fixed localized input has a finite
pole order, so only finitely many terms of the exponential
contribute to its principal part. The output therefore has
algebraic coefficients even for unrestricted \(T\).

On the \(r\)-block summand the internal differential is
\((-1)^r d_{B^{\otimes\ell(\pi)}}\). The total differential
is this map plus the sum of the elementary faces.
It has degree one. To prove that it squares to zero, the
normal contacts must first be compared in a common target.

\subsection{The three normal channels}

Put \(s=z_1-z_3\), \(t=z_2-z_3\) and \(r=s-t\).
For any vector space \(P\), take
\[
 \mathcal G(P)=P[[s,t]][(st(s-t))^{-1}].
\]
Power series are formed before localization, so an element
has a finite pole bound but can have infinitely many
independent state coefficients. The full contact module
\(E_{s,t}(P)\) consists of finite sums of
\(p\,s^{-a}t^{-b}\), \(a,b\geq1\), with their normal
densities and right differential actions.

The change from contact coordinates \((r,t)\) to \((s,t)\)
is
\begin{equation}\label{v4-eq:regular-contact-transfer}
 J(r^{-a}t^{-b})
       =\sum_{j=0}^{b-1}\binom{a+j-1}{j}
                   s^{-a-j}t^{-b+j}.
\end{equation}
Indeed, \(J\) sends the simple contact to \(s^{-1}t^{-1}\),
the derivative in \(r\) to that in \(s\), and the old
derivative in \(t\) to \(\partial_s+\partial_t\).
Applying those derivatives gives the displayed formula.
Define
\[
 \begin{aligned}
 L(F)&=J\operatorname{pp}_t\operatorname{pp}_rF(r+t,t)
                                      &&(r\ll t),\\
 R(F)&=\operatorname{pp}_s\operatorname{pp}_tF(s,t)
                                      &&(t\ll s),\\
 N(F)&=\operatorname{pp}_t\operatorname{pp}_sF(s,t)
                                      &&(s\ll t).
 \end{aligned}
\]
The first principal part in each expression is the rightmost
one. These are the channels \(12\), \(23\) and \(13\),
expressed in one contact module.

\begin{lemma}\label{v4-lem:regular-three-channel}
For every \(F\in\mathcal G(P)\),
\[
                         L(F)-R(F)+N(F)=0.
\]
\end{lemma}

\begin{proof}
Write \(F=G/(s^at^b(s-t)^c)\). A monomial surviving two
principal parts has total degree at most \(-2\).
Consequently numerator terms of degree greater than
\(a+b+c-2\) contribute to none of the three maps.
Only a finite polynomial part of \(G\) matters.

Over \(k[t,t^{-1}]\), the polynomials \(s\) and \(s-t\)
are coprime. Partial fractions reduce to polynomials in
\(s\), terms \(a_j(t)s^{-j}\), and terms
\(b_j(t)(s-t)^{-j}\), with Laurent coefficients in \(t\).
A polynomial in \(s\) gives zero. For \(a_j(t)s^{-j}\),
the expansion at \(r=0\) is regular, so \(L=0\) and
\(R=N=s^{-j}\operatorname{pp}_t(a_j)\).

For \(b_j(t)(s-t)^{-j}\), the expansion in \(s\ll t\)
is regular in \(s\), so \(N=0\). In the other chamber,
\[
 (s-t)^{-j}
       =\sum_{n\geq0}\binom{j+n-1}{n}s^{-j-n}t^n.
\]
A term \(t^{-\beta}\) of \(b_j\), with \(\beta\geq1\),
therefore gives
\(\sum_{n<\beta}\binom{j+n-1}{n}s^{-j-n}t^{n-\beta}\)
in \(R\). This is exactly \(J(r^{-j}t^{-\beta})\),
its contribution to \(L\). The identity follows.
\end{proof}

For \(F=1/(s(s-t))\), the three values are
\(s^{-1}t^{-1},0,-s^{-1}t^{-1}\).
Deleting the outer channel leaves a nonzero defect.
For \(F=1/((s-t)^2t^2)\), they are
\[
 L(F)=R(F)=s^{-2}t^{-2}+2s^{-3}t^{-1},
 \qquad N(F)=0.
\]
The second term comes from the normal-coordinate transfer;
keeping only the scalar residue would lose it.

\begin{theorem}[The regular residue complex]
\label{v4-thm:regular-residue-complex}
The total differential on
\eqref{v4-eq:regular-run-complex} is differential-linear
and squares to zero. Unital differential algebra maps
commuting with \(T\) induce chain maps by their tensor
powers on the runs.
\end{theorem}

\begin{proof}
The derivative of \(e^{uT_{C_i}}\) is
\(T_{C_i}e^{uT_{C_i}}\). Together with
\eqref{v4-eq:regular-run-connection} and
Lemma~\ref{v4-lem:regular-run-functor}, this identifies the
input and output tangential connections in every face.
In the normal direction,
\(\operatorname{pp}_u(-\partial_u f)\) is right
multiplication by \(\partial_u\); differentiation of the
Taylor factor gives precisely the input translation term.
Multiplication by \(u\) agrees with the lowering action
on the \(e_j\). Thus the face extends to a
differential-module map.

The internal square is zero. Each unshifted face is a chain
map, and the two adjacent block numbers have opposite
internal signs. The mixed terms therefore cancel.
Two successive mergers either join disjoint pairs or join
three blocks. For a final partition \(\sigma\), translate
each original run from its original block coordinate to
the retained coordinate of its \(\sigma\)-block. Denote
this formal Taylor operator by \(E_{\pi/\sigma}\).
The derivation identity gives, for every intermediate \(\rho\),
\[
 p_{\rho\sigma}E_{\rho/\sigma}\,
 p_{\pi\rho}E_{\pi/\rho}
       =p_{\pi\sigma}E_{\pi/\sigma}.
\]
Intermediate products only parenthesize the same factors,
and their translation displacements add. On a fixed
localized input, each compared contact uses finitely many
Taylor coefficients, so this formal identity determines
the actual normal distributions.

The common coefficient can now be taken through every
route. For disjoint pairs the two normal orientations
have opposite signs. For three blocks the oriented sum
is the identity of Lemma~\ref{v4-lem:regular-three-channel},
with any remaining coefficients as spectators. These
two cases exhaust the square of the boundary.
Maps of differential algebras preserve every product,
Taylor coefficient and internal differential in these
formulas, proving the final assertion.
\end{proof}

\section{Collision coherence and framed recognition}
\label{v4-sec:regular-recognition}

The complexes just constructed live on the actual position
charts. To obtain an object of \(\mathcal W_X\), their
exceptional collision comparisons must also be supplied.

\begin{theorem}\label{v4-thm:regular-word-realization}
The complexes \(C_I^{\mathrm{run}}(B)\) have canonical
coherent equivalences
\[
 d_\alpha^!C_I^{\mathrm{run}}(B)
       \simeq C_J^{\mathrm{run}}(B)
\]
for every monotone surjection \(\alpha:I\twoheadrightarrow J\).
They form a connected coaugmented cut coalgebra. The
construction is functorial in differential coefficient algebras
and preserves quasi-isomorphisms.
\end{theorem}

\begin{proof}
A partition term survives \(d_\alpha^!\) only when each
\(\alpha\)-fibre lies in one equality block. Otherwise that
collision lands in a diagonal removed by \(j_\pi\), and
the exceptional pullback vanishes by \(i^!j_*=0\).
A surviving partition is uniquely
\(\pi=\alpha^{-1}(\bar\pi)\). Since the fibres of \(\alpha\)
are consecutive, deleting their repetitions identifies
the run lists of \(\pi\) and \(\bar\pi\).
The closed diagonal \(i_\pi\) factors through \(d_\alpha\);
the closed-embedding identity \(d_\alpha^!d_{\alpha,*}
\simeq\mathrm{id}\) and open base change identify the
surviving coefficient with its \(\bar\pi\)-term. Its
density and shift are still \(\omega_{X^{\bar\pi}}\)
and \([|\bar\pi|]\). Formula
\eqref{v4-eq:regular-run-face} identifies the surviving
differentials.

Under two collisions, the same partitions survive as
under their composite, and the same adjacent repetitions
are deleted. The transitivity maps for closed direct
image and open base change therefore give all collision
coherences.

For a consecutive cut \(I=A|D\), restrict to the open
where every position in \(A\) differs from every position
in \(D\). A partition with a block crossing the cut has
empty support there. Every surviving partition is
\(\pi_A\sqcup\pi_D\), and its runs are the concatenation
of the two run lists. Thus there is an equivalence
\[
 C_I^{\mathrm{run}}(B)|_{U_{A,D}}
   \simeq
 \bigl(C_A^{\mathrm{run}}(B)\boxtimes
       C_D^{\mathrm{run}}(B)\bigr)|_{U_{A,D}}.
\]
Use the canonical tensor-shift isomorphisms when
identifying the differentials. Explicitly, the map
from \(M[a]\otimes N[b]\) to \((M\otimes N)[a+b]\)
multiplies \(x\otimes y\), with \(x\) of unshifted
degree \(|x|\), by \((-1)^{b|x|}\).
The block densities occur in their concatenated order.

Localization followed by this equivalence, summed over
cuts, defines the coproduct. For three or more output
blocks every parenthesization restricts to the same
disjoint-block open and concatenates the same run
coefficients. This proves coassociativity and its
compatibility with collisions. Empty cuts are identity
maps on the nonempty coefficient and use
\(C_\varnothing=k\), giving the counit and coaugmentation.
More than \(|I|\) positive output blocks are impossible.

Finally filter by the number of equality blocks, using
the terms with at most \(r\) blocks for the \(r\)-th
stage. The differential preserves this finite filtration.
On its associated graded, a quasi-isomorphism of \(B\)'s
acts by finite tensor powers over a field, followed by
derived coefficient functors. It is a quasi-isomorphism
on every term, and hence on the finite total object.
\end{proof}

The four-point partial strata display the coherence before
the general proof is needed again:
\[
 \begin{array}{c|c|c}
 \pi&\text{run pattern}&\text{coarsening from }1|2|3|4\\ \hline
 12|34&AABB& a\otimes b\otimes c\otimes d\mapsto ab\otimes cd\\
 13|24&ABAB& a\otimes b\otimes c\otimes d\mapsto a\otimes b\otimes c\otimes d\\
 14|23&ABBA& a\otimes b\otimes c\otimes d\mapsto a\otimes bc\otimes d.
 \end{array}
\]
Merging either pair first gives the displayed map.
The final collision gives \(abcd\) in every case.
The corresponding normal residue routes still carry
their orientation signs.

The coefficient unit \(k\to B\) supplies a coalgebra map
\(C^{\mathrm{run}}(k)\to C^{\mathrm{run}}(B)\).
On each partition it inserts the coefficient unit in
every run. This source retains nonempty coefficients:
already its singleton is \(\omega_X[1]\).
It is distinct from the isolated empty unit \(K\).
The singleton unit and every tested regular mode are
therefore retained without assigning the empty
coefficient the value \(B\).

Theorem~\ref{v4-thm:bar-cut-duality} now gives an associative
geometric algebra
\[
 A^{\mathrm{run}}(B)
      =\Omega_{\mathrm{cut}}C^{\mathrm{run}}(B),
 \qquad
 B_{\mathrm{cut}}A^{\mathrm{run}}(B)
      \simeq C^{\mathrm{run}}(B).
\]
Its multiplication is cobar concatenation. The field
multiplication of \(B\) is recovered instead from the
framed two-point boundary
\[
 r_B(f;a,b)=\operatorname{pp}_u\bigl(f(e^{uT}a)b\bigr),
 \qquad
 a_{(n)}b=[e_0]\,r_B(u^n;a,b).
\]
The test \(u^{-1}\) recovers \(ab\); all negative-index
modes are retained in the same way. This recovery uses
the full normal module, not evaluation of a distribution
at \(u=0\).

\begin{definition}\label{v4-def:regular-split-presentation}
A \emph{split presentation with run coefficients} consists
of the following data for a complex \(E\) with a commuting
degree-zero operator \(T\):
the literal direct sums
\eqref{v4-eq:regular-run-complex} with \(E\) in every run;
the connections \eqref{v4-eq:regular-run-connection};
a nonzero closed degree-zero singleton vector \(1\) with \(T1=0\);
the saturated-partition collision identifications of
Theorem~\ref{v4-thm:regular-word-realization}, including
their transitivity; and the cut maps given by concatenation.
It also has a differential extending the shifted
internal differential and a binary contact \(r\).
The contact is \emph{regular} when
\([e_0]r(u^n;a,b)=0\) for every \(n\geq0\).
A morphism is the literal tensor power, on each run
coefficient, of one singleton chain map commuting with
\(T\) and \(1\), compatible with all the displayed maps.
Morphisms are not identified up to chain homotopy.
\end{definition}

The definition supplies the carrier and its comparisons.
It does not assume a multiplication on \(E\). Recover
the candidate product by
\(\mu(a,b)=[e_0]r(u^{-1};a,b)\).

\begin{theorem}[Recognition with a specified presentation]
\label{v4-thm:regular-framed-recognition}
A split presentation with regular binary contact is
isomorphic to \(C^{\mathrm{run}}(B)\) as a framed
presentation exactly when the following equations hold:
\(\mu\) is associative with unit \(1\); \(d\) and \(T\)
are derivations of \(\mu\); every elementary face is
\eqref{v4-eq:regular-run-face} for this product; and
the total differential has no further components.
The inverse construction is \(B=(E,d,T,\mu,1)\).
These constructions give an equivalence of the stated
strict categories.
\end{theorem}

\begin{proof}
Every condition holds for the constructed realization.
Conversely the first two conditions make \(E\) a
differential coefficient algebra. Its fields are
\((e^{uT}a)b\). The face equations identify every
normal contact, and the condition on the remaining
components identifies the full differential.
The prescribed splitting identifies the coefficient
modules, their shifts and connections. The prescribed
saturation and cut maps identify the exceptional
comparisons, their coherences and the coproduct.
This gives an isomorphism of all the framed data.

A morphism of presentations recovers its singleton
map. Compatibility with \(r(u^{-1};a,b)\) makes this
map multiplicative, while the other singleton
requirements preserve \(d,T,1\). Conversely such an
algebra map preserves every construction. The two
assignments are inverse on morphisms as well as objects.
\end{proof}

Prescribing the collision comparisons is necessary.
In the scalar presentation, multiply the comparison
for \(\alpha:I\twoheadrightarrow J\) by
\(\lambda^{|I|-|J|}\), where \(\lambda\in k^\times\).
The exponents add under composition and under cut
decomposition, so all their coherences remain valid.
The coefficients and differentials are unchanged.
For \(\lambda\ne1\), no unital singleton map \(k\to k\)
induces a framed isomorphism with the original presentation.

Even the unframed word carrier has coefficients invisible
in arities one and two. Choose distinct \(a,b\in k\).
For a rational point \(p\in X^n\), let \(\delta_p\)
be its differential-module direct image. The system
\[
 F_0=0,\qquad
 F_n=\bigoplus_{\substack{r_1,r_2,r_3\geq1\\
                          r_1+r_2+r_3=n}}
           \delta_{a^{r_1}b^{r_2}a^{r_3}}
\]
has canonical monotone-collision comparisons. Indeed,
the exceptional pullback of a summand vanishes unless
the point is constant on each collision fibre.
Every surviving point is the unique expansion of a
corresponding point of the shorter word, and closed
direct image gives the transitive identifications.
Thus \(F_1=F_2=0\) but \(F_3=\delta_{(a,b,a)}\ne0\).
There are \(\binom{n-1}{2}\) summands for \(n\geq3\).
Binary data alone are not conservative on this carrier.

\subsection{Strict differentials retain higher products}

The recognition theorem is a statement about literal
presentations. For
\[
 B=k1\oplus(ku\xrightarrow{d}kv),\qquad
 |u|=-1,\quad |v|=0,\quad du=v,\quad dv=0,\quad T=0,
\]
give the second summand square-zero multiplication.
The unital maps that multiply \(u,v\) by a scalar
\(\lambda\) are distinct strict algebra maps.
The augmentation \(B\to k\) is a quasi-isomorphism
and has the same composite with every such map.
Consequently these maps all become the identity after
localization. Strict faithfulness does not imply
faithfulness of strict maps into a derived category.
This example does not decide the faithfulness of the
functor from the already localized source.

The cubic algebra in the prologue supplies a different
test. Its strict endomorphism algebra \(E\), with
\(T=0\), is an allowed coefficient algebra. Hence
\(C^{\mathrm{run}}(E)\) is defined with no finite
state truncation. The two-point contact recovers
composition in \(E\), and its coefficient differential
retains the identities
\[
 DU=DV=0,\qquad U^2=DH,\qquad UH+HU=-V.
\]
In particular the single-valued product
\(\langle[U],[U],[U]\rangle=\{-[V]\}\) from
\eqref{v4-prologue:eq:massey} remains encoded by the
product and its primitive. It is nonzero.
Replacing \(E\) by its graded cohomology algebra with
zero differential would replace that product by zero.

There is therefore no zigzag of framed quasi-isomorphisms
between these two presentations: taking the singleton
and its recovered product would give a zigzag of
quasi-isomorphisms of differential graded algebras,
contradicting the computed Massey product.
A higher presentation can express the same information
in a ternary operation. A strict presentation expresses
it in the differential and the chosen primitives.
Neither description justifies discarding it.

\section{An exchange that permits contextual compression}
\label{v4-sec:regular-contextual-exchange}

Retaining every run is one construction. A different one
keeps one coefficient for each occupied position and
allows the intervening coefficients to enter a merger.
Its basic map is
\begin{equation}\label{v4-eq:regular-exchange}
 c_B(a\otimes b)=ab\otimes1+1\otimes ab-a\otimes b.
\end{equation}
This is the \(R_{1,1,1}\) member of the
D\u{a}sc\u{a}lescu--Nichita operators.\footnote{
G.~Massuyeau and F.~F.~Nichita,
\emph{Yang--Baxter operators arising from algebra
structures and the Alexander polynomial of knots},
Theorem~2.3, p.~2;
\href{https://arxiv.org/abs/math/0404010}{arXiv:math/0404010}.}
We shall need its multiplication identities as well as
its braid relation.

\begin{lemma}\label{v4-lem:regular-contextual-exchange}
For every unital differential graded algebra \(B\),
\(c_B\) is a degree-zero chain map, natural in unital
algebra maps, and
\[
 \begin{gathered}
 c_B^2=1,\qquad \mu c_B=\mu,\qquad
 c_B(a\otimes1)=1\otimes a,\qquad
 c_B(1\otimes a)=a\otimes1,\\
 (c_B)_{12}(c_B)_{23}(c_B)_{12}
       =(c_B)_{23}(c_B)_{12}(c_B)_{23},\\
 c_B(1\otimes\mu)
       =(\mu\otimes1)(1\otimes c_B)(c_B\otimes1),\\
 c_B(\mu\otimes1)
       =(1\otimes\mu)(c_B\otimes1)(1\otimes c_B).
 \end{gathered}
\]
\end{lemma}

\begin{proof}
Multiplication and insertion of the closed unit are
degree-zero chain maps. They give chain compatibility
and naturality. The first line follows by substitution.
Both braid words on \(a\otimes b\otimes c\) reduce to
\[
 \begin{split}
 &abc\otimes1\otimes1+1\otimes ab\otimes c
       +a\otimes bc\otimes1+1\otimes1\otimes abc\\
 &\hspace{12mm}
       -ab\otimes c\otimes1-1\otimes a\otimes bc
       -a\otimes b\otimes c.
 \end{split}
\]
The last two identities reduce, respectively, to
\(abc\otimes1+1\otimes abc-a\otimes bc\) and
\(abc\otimes1+1\otimes abc-ab\otimes c\).
Only associativity and the unit have been used.
\end{proof}

Adjacent \(c_B\)'s therefore give an action of the
symmetric group on each \(B^{\otimes n}\).
Multiplication is invariant under these specified
crossings. This assertion does not replace \(c_B\)
by the usual graded flip.

For a coarsening \(\pi\leq\rho\), start with one
\(B\)-factor per \(\pi\)-block. Stably group the
blocks belonging to the same \(\rho\)-block by
adjacent \(c_B\)'s, then multiply each group in order.
Denote the resulting map by
\[
 q_{\pi\rho}:B^{\otimes|\pi|}
                  \longrightarrow B^{\otimes|\rho|}.
\]
Sorting the runs of each \(\pi\)-block together and
multiplying them similarly gives
\(\kappa_\pi:B^{\otimes\ell(\pi)}\to B^{\otimes|\pi|}\).

\begin{proposition}\label{v4-prop:regular-contextual-coherence}
These maps are independent of the adjacent-transposition
expressions used for the prescribed grouping and obey
\[
 q_{\rho\sigma}q_{\pi\rho}=q_{\pi\sigma},
 \qquad
 \kappa_\rho p_{\pi\rho}=q_{\pi\rho}\kappa_\pi.
\]
They are compatible with reordering the blocks by
the specified symmetric-group action.
\end{proposition}

\begin{proof}
The braid, involution and distant-commutation relations
identify all words for a fixed permutation. Crossings
inside a group disappear after multiplication because
\(\mu c_B=\mu\). The two multiplication identities
in Lemma~\ref{v4-lem:regular-contextual-exchange}
move a merger through a crossing.

Apply these identities to two successive grouping
diagrams. Move crossings before the mergers and combine
successive multiplications by associativity. The result
groups the original inputs directly into their final
blocks. This proves the first identity. Applying the
same argument to the original run list proves the
second. Starting with a reordered list instead gives
the stated order compatibility.
\end{proof}

At the first nonconsecutive collision this gives
\begin{equation}\label{v4-eq:regular-contextual-three}
 \kappa_{13|2}(a,b,c)
       =abc\otimes1+a\otimes bc-ab\otimes c,
 \qquad
 \mu\kappa_{13|2}(a,b,c)=abc.
\end{equation}
For \(a=c=E_{12}\), \(b=E_{21}\), the result is
\[
 E_{12}\otimes1+E_{12}\otimes E_{22}
                  -E_{11}\otimes E_{12}.
\]
Its product is \(E_{12}\). Each term has used the
middle input, so this map does not factor through
the blind outer multiplication \(a\otimes c\mapsto ac\).
The units in the formula are coefficient units at
retained positions; they do not remove a position.

For two interlaced pairs,
\begin{equation}\label{v4-eq:regular-contextual-four}
 q_{13|24}(a,b,c,d)
       =abc\otimes d+a\otimes bcd-ab\otimes cd.
\end{equation}
Its product is \(abcd\). Taking
\(a=c=E_{12}\) and \(b=d=E_{21}\) gives
\[
 2E_{12}\otimes E_{21}-E_{11}\otimes E_{11},
 \qquad \mu q_{13|24}=E_{11},
\]
whereas \(ac\otimes bd=0\).
For the nested pair pattern,
\[
 q_{14|23}(a,b,c,d)
       =abcd\otimes1+a\otimes bcd-abc\otimes d.
\]
Both orders of the two mergers give these formulas,
by the proved grouping identities.

\subsection{The exchange contains more than its product}

For \(B=k[x]/(x^2)\), both \(c_B\) and the usual flip
\(\tau\) have the same composite with multiplication.
In the basis \(1\otimes1,1\otimes x,x\otimes1,x\otimes x\),
\(c_B\) fixes the first vector, exchanges the middle
two and negates the last. Hence
\[
 \operatorname{Tr}c_B=0,\qquad \operatorname{Tr}\tau=2.
\]
The two exchanges are not linearly conjugate.

One can measure the difference from the ordinary flip
for every finite-dimensional ordinary algebra. Put
\(n=\dim B\), \(a_n=n(n-1)/2\), \(q=\dim[B,B]\), and
\[
 s(b)=b\otimes1+1\otimes b,\qquad R_B=\tau c_B=s\mu-\tau.
\]
Here \([B,B]\) is the linear span of all commutators.
Write \(J_m(\lambda)\) for a Jordan block of size \(m\)
with eigenvalue \(\lambda\).

\begin{proposition}\label{v4-prop:regular-exchange-jordan}
The operator \(R_B^2-1\) is square-zero and has rank \(q\).
The Jordan form of \(R_B\) is
\[
 qJ_2(1)\ \oplus\
 (n+a_n-2q)J_1(1)\ \oplus\ a_nJ_1(-1).
\]
Consequently
\[
 \operatorname{rank}(R_B-1)=a_n+q,\qquad
 \operatorname{Tr}R_B=n,\qquad
 R_B^2=1\ \Longleftrightarrow\ B\text{ is commutative}.
\]
\end{proposition}

\begin{proof}
Since \(\tau s=s\) and \(\mu s=2\,\mathrm{id}\),
\[
 R_B^2-1=s(\mu-\mu\tau).
\]
The map \(s\) is injective and
\((\mu-\mu\tau)s=0\). This proves the square-zero
and rank assertions.

On \(\operatorname{Sym}^2B\), the map \(s\mu/2\)
is a projection onto \(sB\). The operator \(R_B\)
is \(1\) on this image and \(-1\) on its kernel,
which has dimension \(a_n\).
On \(\bigwedge^2B\), it is \(1+s\mu\), with
\(s\mu\) taking values in \(sB\) and having rank \(q\).
Thus \(R_B\) restricts to identity plus a square-zero
operator of rank \(q\) on
\(sB\oplus\bigwedge^2B\), a space of dimension \(n+a_n\).
This gives the displayed blocks and all their consequences.
\end{proof}

For \(B=\operatorname{Mat}_d(k)\), commutators span
the trace-zero matrices, so \(q=d^2-1\).
At \(d=2\), the multiplicities of
\(J_2(1),J_1(1),J_1(-1)\) are \(3,4,6\).
These are invariants of the specified coefficient
exchange. Equality after multiplication forgets them.

\section{Descent and the kernel of compression}
\label{v4-sec:regular-descent}

First take \(T=0\). On each word chart of the ordered
atlas \(\mathscr R_I\) of
Section~\ref{v4-sec:word-ordered-atlas}, replace the
run coefficient in \eqref{v4-eq:regular-run-complex}
by \(B^{\otimes|\pi|}\). In its faces replace
\(p_{\pi\rho}\) by \(q_{\pi\rho}\).
Call the resulting chart complex \(C_I^{\mathrm{ctx}}(B)\).
The same three-channel identity proves that its
differential squares to zero: every route to a final
partition has the single coefficient map \(q_{\pi\sigma}\).

On an overlap of two charts, a surviving equality
block has the same internal order in both charts.
Only its order relative to other, separated blocks
can change. Use the action generated by \(c_B\)
to change that block order, together with the
corresponding density and tensor-shift identifications.
These are specified transition maps.

\begin{theorem}\label{v4-thm:regular-contextual-descent}
For \(T=0\), these transitions glue the chart complexes
to a coherent exceptional collision system on
\(\mathscr R_I\). It is a connected coaugmented
coalgebra for the atlas convolution \(\star\).
The maps \(\kappa_\pi\) define a natural map of
cut coalgebras
\begin{equation}\label{v4-eq:regular-compression}
 C^{\mathrm{run}}(B)
       \longrightarrow\operatorname{res}
                              C^{\mathrm{ctx}}(B),
\end{equation}
where restriction uses the colax comparison of
Proposition~\ref{v4-prop:word-atlas-comparison}.
The map is an isomorphism in arities one and two.
\end{theorem}

\begin{proof}
The symmetric-group relations prove the transition
cocycle on every triple overlap. The grouping
identities make transitions commute with all
surviving faces. A target stratum absent from an
overlap has zero restricted face. Hence these are
chain isomorphisms on the full overlaps.

For an ordered-fibre collision, choose an expanded
chart where its fibres are consecutive. The saturated
partitions are precisely the pullbacks of partitions
of the shorter word. Their densities, shifts and
block coefficients agree. The argument in
Theorem~\ref{v4-thm:regular-word-realization} therefore
gives the exceptional comparison there. Compatibility
with transitions makes it a comparison on the actual
atlas collision map, and saturation under composite
collisions proves transitivity.

On an arbitrary cross-disjoint decomposition
\(I=A\sqcup D\), the surviving partitions are
\(\pi_A\sqcup\pi_D\). Regroup their coefficient
blocks using \(c_B\). This is invertible, and
localization into the resulting external product
defines the coproduct. All iterates use the same
final disjoint open and the same regrouping.
The braid and multiplication identities give
coassociativity, collision compatibility and the
empty-unit identities.

The identity
\(\kappa_\rho p_{\pi\rho}=q_{\pi\rho}\kappa_\pi\)
proves that \eqref{v4-eq:regular-compression} is
a chain map and preserves collisions.
On a consecutive cut, sorting the runs splits
into the sorting on its two sides; this proves
compatibility with the cut coproduct.
At one or two positions every run list has
one occurrence of each equality block, so
each \(\kappa_\pi\) is the identity.
\end{proof}

This construction supplies a chosen exchange.
It does not extend a different prescribed overlap
map by leaving it unchanged.

\begin{proposition}[A split compression kernel]
\label{v4-prop:regular-compression-kernel}
The coefficient map \(\kappa=\kappa_{13|2}\) is
a split surjection for every unital differential
graded algebra. A section and the resulting
projection onto its kernel are
\[
 s_0(x\otimes y)=x\otimes y\otimes1,
 \qquad P=1-s_0\kappa.
\]
Explicitly,
\[
 \begin{split}
 P(a,b,c)={}&a\otimes b\otimes c
      -abc\otimes1\otimes1-a\otimes bc\otimes1\\
      &+ab\otimes c\otimes1.
 \end{split}
\]
For an ordinary algebra of finite dimension \(d\),
\(\dim\ker\kappa=d^3-d^2\).
\end{proposition}

\begin{proof}
Substituting \(c=1\) in
\eqref{v4-eq:regular-contextual-three} gives
\(\kappa s_0=\mathrm{id}\).
Thus \(P^2=P\), \(\kappa P=0\), and \(P\) is
identity on \(\ker\kappa\).
All maps commute with the differential.
The dimension follows from the split surjection.
\end{proof}

On the open equality stratum
\(S=\{z_1=z_3\ne z_2\}\), all other partition terms
vanish under exceptional restriction. Both remaining
terms have shift \([2]\). Consequently
\[
 i_S^!\operatorname{Cone}
 \bigl(C_3^{\mathrm{run}}(B)\to
       (\operatorname{res}C^{\mathrm{ctx}}(B))_3\bigr)
       \simeq\omega_S\otimes\ker\kappa[3].
\]
The extra shift is the shift of a kernel in the
cone of a split surjection. For \(d>1\), this
proves that compression is not an equivalence.

\begin{proposition}[An obstruction to unchanged descent]
\label{v4-prop:regular-unchanged-descent}
Let \(B\) be an ordinary finite-dimensional
algebra with \(T=0\) and \(\dim B=d>1\).
The run complexes on the separate word charts
cannot descend to \(\mathscr R_I\) with their
chart objects unchanged.
\end{proposition}

\begin{proof}
The charts \(123\) and \(132\) overlap wherever
\(z_2\ne z_3\), and both contain \(S\).
The run list of \(13|2\) is \(ABA\) in the
first chart and \(AB\) in the second.
Their exceptional restrictions to \(S\) are
\[
 \omega_S[2]\otimes B^{\otimes3},
 \qquad
 \omega_S[2]\otimes B^{\otimes2}.
\]
They have ranks \(d^3\) and \(d^2\) in the
same degree. An overlap equivalence would
induce a quasi-isomorphism between them,
which is impossible.
\end{proof}

For a differential graded algebra the comparison
must use cohomology, not the dimensions of the
underlying complexes. The same argument applies
when \(H^*(B)\) has finite total dimension
greater than one. In contrast, the square-zero
contractible extension of \(k\) considered above
has run complexes quasi-isomorphic to the scalar
ones. Its underlying state dimensions give no
obstruction of this kind.

\section{Translation frames and the algebraic boundary}
\label{v4-sec:regular-translation}

The run construction allowed every derivation \(T\).
For contextual descent, translation also enters the
overlap maps. Suppose first that \(T\) is locally
nilpotent, meaning that for each \(b\in B\) some
power \(T^Nb\) is zero. Then
\[
 E_z=e^{zT}:\mathcal O_X\otimes B
                  \longrightarrow\mathcal O_X\otimes B
\]
is an algebraic change of frame, with inverse
\(E_{-z}\), and
\[
                  \partial_z E_z=E_z(\partial_z+T).
\]
This is an identity of operators on sections.
It changes the translation connection to the
constant connection.

Conversely, an algebraic horizontal frame \(E(z)\)
with \(E(0)=1\) is necessarily this one.
Writing \(E(z)=\sum_n z^nE_n\) on each input,
horizontality gives
\((n+1)E_{n+1}=E_nT\), hence \(E_n=T^n/n!\).
Polynomiality on every vector is exactly local
nilpotence. Thus the hypothesis is precise for
this method of constructing an algebraic frame.

On a partition put
\[
 E_{\pi,\mathrm{run}}
       =\bigotimes_{r=1}^{\ell(\pi)}
                         e^{z_{\beta_\pi(r)}T},
 \qquad
 E_{\pi,\mathrm{block}}
       =\bigotimes_{C\in\pi}e^{z_CT}.
\]
Conjugate every contextual face and overlap map
by its source and target frames, and set
\begin{equation}\label{v4-eq:regular-translated-compression}
 \kappa_\pi^T
       =E_{\pi,\mathrm{block}}^{-1}
                         \kappa_\pi E_{\pi,\mathrm{run}}.
\end{equation}

\begin{theorem}\label{v4-thm:regular-translated-descent}
For locally nilpotent \(T\), the contextual chart
complexes with these conjugated maps descend to
the ordered atlas and form a connected
\(\star\)-coalgebra. Formula
\eqref{v4-eq:regular-translated-compression}
gives a map of cut coalgebras from
\(C^{\mathrm{run}}(B)\) to its restriction.
It retains every regular field coefficient.
\end{theorem}

\begin{proof}
All frames are algebraic and invertible on
sections and commute with \(d\).
At a merged position,
\[
                  e^{zT}\mu=\mu(e^{zT}\otimes e^{zT}).
\]
Conjugating a face therefore gives exactly its
normal Taylor transport. Intermediate frames
cancel in every composition identity, including
the transition cocycle, the grouping equations
and the collision and coproduct comparisons.
Thus Theorem~\ref{v4-thm:regular-contextual-descent}
applies in the translation frames.
At two positions the resulting contact is
\(\operatorname{pp}_u(f(e^{uT}a)b)\), so its
Laurent tests recover all the original modes.
\end{proof}

For two positions \(t,0\), written in the exchanged
order \(0,t\), the transition is
\begin{equation}\label{v4-eq:regular-spectral-exchange}
 \begin{split}
 c_T(t)(a\otimes b)
       ={}&(e^{tT}a)b\otimes1
           +1\otimes a(e^{-tT}b)\\
         &-(e^{tT}a)\otimes(e^{-tT}b).
 \end{split}
\end{equation}
The reverse transition uses \(-t\). Conjugating
the constant relations proves
\[
 c_T(-t)c_T(t)=1,
 \qquad
 c_{12}(v)c_{23}(u+v)c_{12}(u)
       =c_{23}(u)c_{12}(u+v)c_{23}(v).
\]
These identities hold on all states; they are
the cancellation of intermediate Taylor frames
around the two placed diagrams.

At the partial collision with coordinates
\((t,0,t)\), formula
\eqref{v4-eq:regular-translated-compression} becomes
\begin{equation}\label{v4-eq:regular-translated-three}
 \begin{split}
 \kappa_{13|2}^T(a,b,c;t)
       ={}&a(e^{-tT}b)c\otimes1
                +a\otimes b(e^{tT}c)\\
          &-a(e^{-tT}b)\otimes e^{tT}c.
 \end{split}
\end{equation}
The binary evaluation
\(x\otimes y\mapsto(e^{tT}x)y\)
gives \((e^{tT}a)b(e^{tT}c)\):
the second and third terms cancel, and the first
gives that ordered product. Local nilpotence
makes the displayed coefficient algebraic.

\begin{example}[A nonconstant polynomial exchange]
\label{v4-ex:regular-polynomial-exchange}
Take \(B=k\langle x,y\rangle\), \(Tx=0\), \(Ty=x\).
Translation lowers the number of \(y\)'s in a word,
so it is locally nilpotent. Here
\(Y_B(y,t)y=y^2+txy\), and
\[
 \begin{split}
 c_T(t)(y\otimes y)
       ={}&y^2\otimes1+1\otimes y^2-y\otimes y\\
          &+t(xy\otimes1-1\otimes yx
                         -x\otimes y+y\otimes x)
             +t^2x\otimes x.
 \end{split}
\]
This is a polynomial transition on the original
algebraic overlap, with the proved braid and
reverse-transition identities. Its position
dependence is given by the stated Taylor frames;
no nontrivial full-turn monodromy follows from it.
\end{example}

For a derivation that is not locally nilpotent,
the same formula can instead be only formal.
For example, on \(B=\operatorname{Mat}_2(k)\)
take \(T=\operatorname{ad}\operatorname{diag}(1,0)\).
Then
\[
 c_T(t)(E_{12}\otimes E_{21})
       =e^tE_{11}\otimes1+e^t1\otimes E_{11}
                       -e^{2t}E_{12}\otimes E_{21}.
\]
Its formal coefficients satisfy the same identities.
They are not algebraic functions on
\(\mathbb A^2\setminus\Delta\). The run construction
still exists, because each of its normal faces
uses only a finite Taylor jet.

\subsection{A compression excluded by unrestricted translation}

An \emph{algebraic two-state expression} for an output
at positions \(t,0\), rooted at zero, is a finite sum
with coefficients in \(k((t))\) of
\[
 E_t^{\mathrm L}(x,y)=(e^{tT}x)y,\qquad
 E_t^{\mathrm R}(x,y)=x(e^{tT}y),
\]
where the pairs \(x,y\in B\) are fixed states.
A coefficientwise series in
\((B\otimes B)[[t]]\) is a different permitted input.

\begin{theorem}\label{v4-thm:regular-unrestricted-obstruction}
Let
\[
 B=k\langle a_i,b_i,c_i\mid i\geq0\rangle,\qquad
 Ta_i=a_{i+1},\quad Tb_i=b_{i+1},\quad Tc_i=c_{i+1}.
\]
The ordered output
\[
                      F(t)=(e^{tT}a_0)b_0(e^{tT}c_0)
\]
has no algebraic two-state expression.
\end{theorem}

\begin{proof}
Project to the span of words \(a_i b_j c_k\) and
identify that span with \(k[p,q,r]\) by
\(a_i b_j c_k\mapsto p^iq^jr^k\).
The output becomes \(e^{t(p+r)}\).
In a binary expression, the translated state
is a contiguous prefix or suffix of the output
word. Its projection is therefore a finite sum
of terms
\[
 f(t)P(p,q,r)e^{tL},\qquad
 L\in\{0,p,p+q,p+q+r,r,q+r\},
\]
with polynomial \(P\) and Laurent \(f\).
This includes empty factors and arbitrary
descendants in the two fixed states.
Other generator patterns contribute zero.

Set \(D=\partial_p+2\partial_q+4\partial_r\).
The six exponential eigenvalues are
\(st\), \(s\in S=\{0,1,3,4,6,7\}\).
If all the finitely many polynomials \(P\)
have degree at most \(d\), the commuting operator
\[
                       \prod_{s\in S}(D-st)^{d+1}
\]
kills the proposed expression. Applied to
\(e^{t(p+r)}\), it instead gives
\[
          (80t^6)^{d+1}e^{t(p+r)}\ne0,
\]
since \(D\) has eigenvalue \(5t\) there.
These calculations hold in
\(k((t))[[p,q,r]]\). They contradict the proposed
finite expression.
\end{proof}

The obstruction fixes the binary coefficient
class in advance. It does not exclude every
compression for a non-nilpotent translation.
For a commutative coefficient algebra, multiplying
the runs of each block already gives horizontal
block coefficients, using the ordinary flip.
The derivation and commutativity identities make
these maps compatible with all coarsenings.
For instance, \(B=k[x]\), \(Tx=x\), has no
polynomial horizontal frame, yet
\[
           (e^{tT}a)b(e^{tT}c)=(e^{tT}(ac))b
\]
is a two-state expression.

For general \(B\), allowing the formal expression
\eqref{v4-eq:regular-translated-three} changes the
coefficient tensor and its completion. Descent then
requires actual continuous collision functors and
overlap maps on that coefficient category, together
with a comparison preserving all the field modes,
vacuum data and cut maps. The algebraic obstruction
above leaves that completed comparison open.
The regular run construction itself needs none of
those additional choices.

The next construction introduces a specified singular
local field factor. Its poles alter the normal
principal parts. The ordered coefficient words,
their run maps and the distinctions between the
empty unit, the local vacuum and a prescribed
state remain necessary.
